\begin{center}\begin{tikzpicture}[scale=.73,font=\scriptsize]
\shade[inner color=newsblue!6,outer color=newsblue!20] (0,0) circle (2.05);
\draw[newsblue!70,thick] (0,0) circle (2.05);
\draw[red!65!black,dashed,very thick] (0,0) circle (1.35);
\fill[blue!45!green!35] (0,0) circle (.73);
\fill[newsblue!80] (0,0) circle (.38);
\node[white] at (0,0) {地球};
\draw[->,very thick,green!55!black] (1.72,1.1) .. controls (1.35,.5) and (1.45,-.4) .. (.96,-.85);
\node[green!40!black,align=left,anchor=west] at (2.2,1.25) {合唱波与局部加速\\（概念示意）};
\node[red!65!black,anchor=west] at (2.1,-.45) {经验性屏障};
\draw[->,red!65!black] (2.0,-.45)--(1.26,-.45);
\node[gray,align=center,text width=6.8cm] at (0,-2.8) {非比例轨道图；电子可在边界内侧被加速，并非穿透实体墙。依据：京都大学2026年10月27日研究说明。};
\end{tikzpicture}\end{center}